% !TeX root = ../main.tex
\section{Stride and Gather Access}\label{sec:stride}

\begin{hypothesis}
  Unit stride lets each \qty{64}{\byte} cache line serve \num{16} consecutive float32 elements, so
  SIMD helps. As the stride grows, each vector still needs \num{8} elements but they come from
  further apart: once the stride reaches \num{16} elements (\qty{64}{\byte}) every element sits in its
  own cache line and effective bandwidth drops by up to $16\times$. We therefore expect time per
  processed element to rise roughly linearly with stride until it saturates at one cache line per
  element, and we expect the SIMD advantage to vanish for any non-unit stride, because the compiler
  must gather lanes individually. Gather (random permutation) should be the worst case.
\end{hypothesis}

\begin{knobs}
  Fixed       & SAXPY; float32; native and scalar builds \\
  Varied      & access (unit, stride $1/2/4/8/16$, gather); $N$ cache-resident (\num{32768}) and
                DRAM-resident (\num{4194304}) \\
  Definition  & stride $s$ processes elements $0, s, 2s, \ldots$ over an array of $N s$ elements;
                gather applies SAXPY through a random permutation index so every element is touched
                once with no spatial locality. \\
  Script      & \code{SIMD_Advantage_Profiling/scripts/run.sh stride} \\
\end{knobs}

\begin{figure}[tbp]
  \centering
  \IfDataExists{stride}{%
    \begin{tikzpicture}
      \begin{axis}[acs, xmode=log, log basis x=2, xlabel={Stride (elements)}, ylabel={Time per processed element (ns)}, legend pos=north west]
        \addplot+[iqr bars] table[x=stride, y=ns_per_elem_median, iqr of=ns_per_elem, only if={access}{stride}, only if={variant}{scalar}, only if={n}{4194304}] {data/stride.csv};
        \addlegendentry{scalar (DRAM)}
        \addplot+[iqr bars] table[x=stride, y=ns_per_elem_median, iqr of=ns_per_elem, only if={access}{stride}, only if={variant}{native}, only if={n}{4194304}] {data/stride.csv};
        \addlegendentry{native (DRAM)}
      \end{axis}
    \end{tikzpicture}%
  }
  \caption{Effect of stride on native and scalar SAXPY time per \emph{processed} element, DRAM-resident
  ($N=\num{4194304}$). The two curves coincide for every non-unit stride.}\label{fig:stride}
\end{figure}

\begin{table}[tbp]
  \centering\small
  \caption{Time per processed element (ns), DRAM-resident SAXPY, $N=\num{4194304}$.}\label{tab:stride}
  \begin{tabular}{@{}lrrr@{}}
    \toprule
    Access & scalar & native & native speedup \\
    \midrule
    unit (stride 1)   & 0.453 & 0.252 & $1.80\times$ \\
    stride 2          & 0.790 & 0.787 & $1.00\times$ \\
    stride 4          & 1.540 & 1.509 & $1.02\times$ \\
    stride 8          & 2.944 & 2.977 & $0.99\times$ \\
    stride 16         & 6.074 & 6.122 & $0.99\times$ \\
    gather (random)   & 10.83 & 10.69 & $1.01\times$ \\
    \bottomrule
  \end{tabular}
\end{table}

\subsection{Analysis}
\Cref{tab:stride} makes the pattern stark. Only unit stride lets \code{native} beat \code{scalar}
(\,$1.80\times$\,); at stride~2 and beyond the two builds are within measurement noise of each other.
Two things happen at once. First, the effective bandwidth collapses: time per processed element
roughly doubles from stride~2 to~4 to~8 to~16 (\num{0.79} $\to$ \num{1.51} $\to$ \num{2.98}
$\to$ \qty{6.12}{\nano\second}), because once the stride reaches \num{16} elements each access pulls a
fresh \qty{64}{\byte} line and $15/16$ of every line is wasted. Second, the SIMD advantage
disappears: a strided or gathered load cannot be a single contiguous \code{vmovups}, so the compiler
either falls back to scalar element loads or (for gather) emits per-lane address arithmetic; the
vector ALU is then starved regardless of width. The disassembly confirms this---\code{native}'s
gather kernel stays on \qty{128}{bit} \code{xmm} and never reaches \code{ymm}
(\cref{tab:vec-summary}), and GCC does not emit an AVX2 \code{vgather} here at all. Gather is the
worst case at \qty{10.7}{\nano\second} per element, $\approx\!42\times$ slower than unit stride,
because the random permutation defeats the hardware prefetcher as well. The practical lesson is that
access pattern, not vector width, is the first-order knob for these kernels.
